\scanpage{51}
Soit donc un ``triangle'' subordonné à
$(K,(D_i)_{i\in I})$, \deleted{\unsure{$f$}} défini par $J\in\unsure{\mathfrak P_3|\mathcal P_3}(I)$, formé
de trois ps. droites concourantes. On
veut décrire \deleted{son} l'orientation ``canonique''
de $K_J$.

\deleted{\unsure{1er)}} 1\textsuperscript{er} cas) Il existe une \deleted{$i\in I$} ps. droite
\deleted{\unsure{$\Theta$}} $\Delta$, qui ne passe par aucun des
trois sommets de $K_J$ — i.e. telle que
le système $K_{J\cup\{i\}}$ de 4 ps. droites
soit formé d'un syst. de \unsure{qua.|quatre} ps. droites
3 à 3 non concourantes. On utilise

\uline{Proposition} Le \deleted{catégorie des arrangem\textsuperscript{ts}} \inserted{foncteur, \deleted{\unsure{qui à la}}\inserted{qui à} qui précède}
\deleted{1-\unsure{dim}} \deleted{gén. de 4 ps. droites}
\deleted{\unsure{ne pas}}
\[\left(\begin{gathered}\text{arrangements simples}\\\text{de 4 ps. droites}\end{gathered}\right)
\longrightarrow
\left(\begin{gathered}\text{arr. 1-dimensionnels}\\\text{simples de 4}\\\text{pss-droites}\end{gathered}\right)\]
est une équivalence de catégories.

(On le démontre en comparant les
deux \deleted{\unsure{catégories}}\unsure{groupoïdes}, à celui des ens. \deleted{finis} \inserted{à}
4 éléments)

Donc, comme $K_{J\cup\{i\}}$ provient canoniquement
d'un arrangem\textsuperscript{t} 2-dimensionnel, $K_J$ hérite
d'une orientation, soit $\omega_{\deleted{\unsure{i}}J,i}$, et on aura
% archive sheet number: 50
